\documentclass{beamer}

% Choix du thème
\usetheme{Madrid}

% Packages utiles
\usepackage[utf8]{inputenc}
\usepackage[T1]{fontenc}
\usepackage[french]{babel}
\usepackage{graphicx} % Pour inclure des images
\usepackage{transparent}

% Titre et informations principales
\title{Séquences de de Bruijn}
\subtitle{Pour la cryptographie}
\author{Pegliasco, Haouas, Marçais, Nguyen}
\date{\today}

\begin{document}

% Page de titre
\begin{frame}
    \titlepage
\end{frame}

% Table des matières
\begin{frame}{Plan}
    \tableofcontents
\end{frame}

\section{Introduction}
\begin{frame}{Introduction}
    \begin{itemize}
        \item Le graphe
        \item La séquence
        \item Attaque brute force
        \item stream cipher
        \item Transformation des données (Blocs et mapping)
        \item Substitution et chiffrement
    \end{itemize}
\end{frame}

\section{Le graphe de de Bruijn}
\begin{frame}{Le graphe de de Bruijn}
    \begin{figure}
        \centering
        \includegraphics[width=0.7\linewidth]{graph.png}
        \caption{Graphe de de Bruijn $B(2, 3)$.}
        \label{fig:graph}
    \end{figure}
\end{frame}

\section{Séquence de de Bruijn}
\begin{frame}{Séquence de de Bruijn}
    \begin{figure}
        \centering
        \includegraphics[width=0.5\linewidth]{sequence.png}
        \caption{Séquence de de Bruijn à partir du graphe $B(2, 2)$}
        \label{fig:sequence}
    \end{figure}
\end{frame}


\section{Attaque brute force}
\begin{frame}{Attaque brute force}
    \begin{figure}
        \centering
        \includegraphics[width=0.7\linewidth]{keypad.jpg}
        \caption{Attaque brute force}
        \label{fig:brute_force}
    \end{figure}
\end{frame}

\section{Chiffrement par flux}
\begin{frame}{Concept général}
    \begin{figure}
        \centering
        \includegraphics[width=0.7\linewidth]{cipher_1.png} 
        \caption{Concept géneral du chiffrement par flux}
        \label{fig:cipher_1}
    \end{figure}
\end{frame}

\begin{frame}{Seed issue de la séquence}
    \begin{figure}
        \centering
        \includegraphics[width=0.7\linewidth]{cipher_2.jpg} 
        \caption{Seed issue de la séquence}
        \label{fig:cipher_2}
    \end{figure}
\end{frame}





\section{Mappage bloc -> indice}
\begin{frame}{Mappage bloc -> indice}
    \begin{figure}
        \centering
        \includegraphics[width=0.5\linewidth]{peg_1.jpeg}
        \caption{Mappage bloc -> indice}
        \label{fig:peg_1}
    \end{figure}
\end{frame}

\section{Substitution et chiffrement}
\begin{frame}{Substitution et chiffrement}
    \begin{figure}
        \centering
        \includegraphics[width=0.5\linewidth]{peg_2.jpeg}
        \caption{Substitution et chiffrement}
        \label{fig:peg_2}
    \end{figure}
\end{frame}


% Dernière diapositive
\begin{frame}{}
    \centering
    Merci pour votre attention !
\end{frame}

\end{document}
